\section{Reward}\label{sec:reward}

At the end of each epoch the epoch budget is divided among Bonfires according to how widely the constructions they
originated are held by other Bonfires, and each Bonfire divides what it receives between the nodes that serve it
and its stakeholders.

\paragraph{Breadth.} The \emph{breadth} of a construction in an epoch is the number of Bonfires that hold it
(Section~\ref{sec:adoption}), counting only Bonfires that are independent of its originator and of one another and
that carry at least a minimum bond (Section~\ref{sec:bond}). A Bonfire counts once in an epoch however often it
used the construction, and it counts again in every later epoch in which it still holds it. Use inside the
originating Bonfire counts nothing: a community needs no payment for work it does for itself.

\paragraph{Credit.} A construction's \emph{credit} in an epoch depends only on its breadth. Below a threshold
breadth it is nothing, and above the threshold it rises with breadth. Writing $n$ for the breadth, $n_0$ for the
threshold and $g$ for the rule,
\begin{equation}\label{eq:credit}
  \text{credit} \;=\; g(n), \qquad g(n) = 0 \ \text{ for } n < n_0, \qquad g \text{ increasing for } n \ge n_0 .
\end{equation}
The threshold is what makes a small group's agreement worth nothing. Two Bonfires that adopt each other's
constructions, or one Bonfire that registers ten thousand constructions which a single partner holds, sit below it.
The shape of $g$ above the threshold is left open (Section~\ref{sec:open}). A share $\lambda$ of a construction's
credit passes to the parents it declares (Section~\ref{sec:lineage}), and a Bonfire's credit is the sum over the
constructions it originated, after those transfers.

\paragraph{Why holding, and not spread.} A construction is paid for lasting in other communities, and nothing is
paid for passing it on. That a convention is taken up does not show that it carries knowledge: among the early
users of a new term, short-lived fashions and lasting changes look alike \citep{tarrade2024position}. A construction that is tried and dropped stops counting when it is dropped. One that a
community keeps is counted afresh each epoch, so a grammar that stays in use keeps being funded. Counting each
Bonfire once also makes repetition worthless. A finding about words points the same way: over short periods, a
word's spread across users and threads predicts its fate better than its frequency does
\citep{altmann2011niche}.

\paragraph{Failure.} A failed reading counts for nothing, and it lowers the standing of the constructions tried
(Section~\ref{sec:constructions}), so a construction pushed into contexts where it does not work leaves that
inventory. Failures are never subtracted from an originator's credit. If they were, a Bonfire could harm another
by adopting its constructions and misusing them.

\paragraph{Division.} A fixed fraction of the epoch budget is set aside to pay for spot checks
(Section~\ref{sec:bond}). The rest is divided among Bonfires in proportion to credit, and each Bonfire's payment is
capped by the ceiling of Section~\ref{sec:bond}. The Bonfire passes a share $s$ of its payment to the nodes pledged
to it, equally, and the rest goes to its stakeholders. The Bonfire publishes $s$. Bonfires compete for nodes
through this rate, so the split between those who structure and those who serve is set by the participants.

\paragraph{From contribution to stake.} A registration names the members whose contributions yielded the
construction: those who taught it, and those who confirmed it. Of the stakeholders' part of what the construction earns, a fraction published by the Bonfire vests to them
as stake in that Bonfire. A member who arrives with nothing and whose words become
constructions that other communities keep thereby becomes a stakeholder. This is the protocol's answer to how
people who take part in making sense of something come to hold a share of what is made.

\paragraph{A worked epoch.} As a thought experiment, let $n_0 = 3$ and, for illustration only, $g(n) = n - 2$ at
or above the threshold. Let $\lambda = 0.1$, the ceiling rate $c = 0.01$ and each node's pledge $q = 10{,}000$ (Section~\ref{sec:bond}), and $s = 0.4$, with $1{,}000$ left of the epoch budget after spot checks. Bonfire $A$ originated $x$, held by six
Bonfires, and $y$, held by two; five nodes are pledged to it. Bonfire $B$ originated $z$, a variant of $x$ held by
four Bonfires; four nodes are pledged to it. Table~\ref{tab:epoch} follows the division.

\begin{table}[!htb]
\centering\small
\begin{tabular}{@{}lrr@{}}
\toprule
 & \textbf{Bonfire $A$} & \textbf{Bonfire $B$} \\
\midrule
Credit from breadth & $x$: $4$, \ $y$: $0$ & $z$: $2$ \\
After the lineage share ($0.2$ from $z$ to $x$) & $4.2$ & $1.8$ \\
Share of $1{,}000$ in proportion to credit & $700$ & $300$ \\
Ceiling, $c\,k\,q$ & $500$ & $400$ \\
Paid & $500$ & $300$ \\
\quad to nodes (share $s$) & $200$ & $120$ \\
\quad to stakeholders & $300$ & $180$ \\
\bottomrule
\end{tabular}
\caption{A worked epoch (thought experiment). The $200$ that Bonfire $A$'s ceiling does not admit is not paid;
what becomes of it is a question of issuance, left to the token model. Construction $y$, held by two Bonfires,
earns nothing.}\label{tab:epoch}
\end{table}
